\begin{proof}

\begin{enumerate}[label={$\langle1\rangle$\arabic*.},leftmargin=*]
\item We first verify that a Hopf surface $H$ with $a(H)=1$ satisfies
$\kappa(H,K_H^{-1})=1$. On a primary Hopf cover the contraction
normal forms are diagonal, or
\[
A(z,w)=(\alpha z+\lambda w^n,\beta w),\qquad
\lambda\ne0,\quad \alpha=\beta^n.
\]
The invariant bivectors $zw\,\partial_z\wedge\partial_w$ in the
diagonal case and $w^{n+1}\,\partial_z\wedge\partial_w$ in the
second case give a nonzero effective anticanonical divisor.
The invariance is an equality under the derivative of the
displayed contraction: its determinant is $\alpha\beta$, and
the displayed coefficient transforms by the same scalar.
These normal forms and the elliptic-fibration description are
also used in \cite[Proposition~3.1 and its proof]{SS2002}.

\item The algebraic reduction of an elliptic Hopf surface is a
holomorphic map to $\PP^1$, with smooth reductions of all fibres.
Every curve is vertical: its intersection with a general fibre
is zero since $H^2(H;\R)=0$, whereas a horizontal curve would
have positive intersection equal to its degree over the base.
Thus the nonzero anticanonical divisor is a positive sum of
reduced fibres. If $m_iD_i$ is the full fibre supported on
$D_i$, each $\OO(m_iD_i)$ is the pullback of $\OO_{\PP^1}(1)$.
A common multiple of the $m_i$ therefore makes a positive power
of $K_H^{-1}$ the pullback of a positive-degree line bundle on
$\PP^1$. Its Iitaka dimension is one.

\item For a secondary Hopf surface take a normal finite primary cover
$p:H'\to H$. The cover is unramified, so
$K_{H'}^{-1}=p^*K_H^{-1}$, and $a(H')=a(H)=1$ by invariance of
transcendence degree under a finite map. The anticanonical
section ring of $H$ is the invariant subring of that of $H'$.
The latter is integral over the invariant subring: a homogeneous
section $s$ satisfies the monic orbit polynomial
$\prod_{g\in\operatorname{Gal}(H'/H)}(T-gs)$. Their fraction
fields have the same transcendence degree, proving
$\kappa(H,K_H^{-1})=1$ also in this case.

\item Now, for each positive integer $m$, the locus
\[
Z_m=\{c\in C^\circ:
 h^0(Y_c,K_Y^{-m}|_{Y_c})\ge2\}
\]
is a closed analytic subset of the smooth base $C^\circ$.
If the very general Hopf fibre had algebraic dimension one,
the preceding calculation would make the union of these
countably many loci contain the complement of a countable
subset of $C^\circ$. A proper closed analytic subset of a
curve is discrete and countable, so some $Z_m$ is the whole
connected smooth base. Coherent base change then gives
$\operatorname{rank}f_*K_Y^{-m}\ge2$.
The relative linear system defines a meromorphic map of $Y$
to $\PP(f_*K_Y^{-m})$ over $C$ with positive-dimensional image
on a general fibre. Its image has dimension at least two.
The projective bundle over the compact curve is projective,
so pulling back its meromorphic functions gives $a(Y)\ge2$,
contrary to the hypothesis. In fact the same argument shows
that, when the smooth fibres are Hopf surfaces, the values with
$a(Y_c)=1$ form at most a countable set.

\end{enumerate}
\end{proof}
